
\section*{2. Polynômes de Snapper}

\subsection*{2.1.}
Rappelons d'abord la définition et quelques propriétés élémentaires des applications polynômes :

\textbf{Définition 2.1.1.} \uline{Soit} @@M0@@ \uline{une application d'un groupe abélien} @@M1@@ \uline{dans un groupe abélien} @@M2@@\uline{. On dit que} @@M3@@ \uline{est une} application polynôme\uline{, s'il existe, pour toute suite} @@M4@@ (@@M5@@\uline{) d'éléments de} @@M6@@\uline{, une famille} @@M7@@ \uline{d'éléments de} @@M8@@\uline{, presque tous nuls, et satisfaisant à la condition suivante : Pour toute suite} @@M9@@ \uline{d'entiers, on a}
@@M10@@
\uline{On dit que} @@M11@@ \uline{est de degré} @@M12@@\uline{, si quels que soient} @@M13@@ \uline{et} @@M14@@\uline{, on a}
@@M15@@
\uline{pour} @@M16@@.

\textbf{Lemme 2.1.2.} \uline{Tout homomorphisme de groupes abéliens est une application polynôme de degré} @@M17@@\uline{. Si} @@M18@@\uline{, où} @@M19@@ \uline{est une application polynôme de degré} @@M20@@ \uline{et} @@M21@@ \uline{une application polynôme de degré} @@M22@@\uline{, alors} @@M23@@ \uline{est une application polynôme de degré} @@M24@@.

Démonstration triviale.

\textbf{Lemme 2.1.3.} \uline{Soit}
@@M25@@
\uline{une suite exacte de groupes abéliens telle que} @@M26@@ \uline{est divisible par}
